\section{The alternating {\v C}ech complex}
\label{section-alternating-cech}

\noindent
This section compares the {\v C}ech complex with the alternating {\v C}ech
complex and some related complexes.

\medskip\noindent
Let $X$ be a topological space. Let $\mathcal{U} : U = \bigcup_{i \in I} U_i$
be an open covering. For $p \geq 0$ set
$$
\check{\mathcal{C}}_{alt}^p(\mathcal{U}, \mathcal{F})
=
\left\{
\begin{matrix}
s \in  \check{\mathcal{C}}^p(\mathcal{U}, \mathcal{F})
\text{ such that }
s_{i_0 \ldots i_p} = 0 \text{ if } i_n = i_m \text{ for some } n \not = m\\
\text{ and }
s_{i_0\ldots i_n \ldots i_m \ldots i_p}
=
-s_{i_0\ldots i_m \ldots i_n \ldots i_p}
\text{ in any case.}
\end{matrix}
\right\}
$$
We omit the verification that the differential $d$ of
Equation (\ref{equation-d-cech}) maps
$\check{\mathcal{C}}^p_{alt}(\mathcal{U}, \mathcal{F})$ into
$\check{\mathcal{C}}^{p + 1}_{alt}(\mathcal{U}, \mathcal{F})$.

\begin{definition}
\label{definition-alternating-cech-complex}
Let $X$ be a topological space. Let $\mathcal{U} : U = \bigcup_{i \in I} U_i$
be an open covering. Let $\mathcal{F}$ be an abelian presheaf on $X$.
The complex $\check{\mathcal{C}}_{alt}^\bullet(\mathcal{U}, \mathcal{F})$
is the {\it alternating {\v C}ech complex} associated to $\mathcal{F}$ and the
open covering $\mathcal{U}$.
\end{definition}

\noindent
Hence there is a canonical morphism of complexes
$$
\check{\mathcal{C}}_{alt}^\bullet(\mathcal{U}, \mathcal{F})
\longrightarrow
\check{\mathcal{C}}^\bullet(\mathcal{U}, \mathcal{F})
$$
namely the inclusion of the alternating {\v C}ech complex into the
usual {\v C}ech complex.

\medskip\noindent
Suppose our covering $\mathcal{U} : U = \bigcup_{i \in I} U_i$ comes
equipped with a total ordering $<$ on $I$. In this case, set
$$
\check{\mathcal{C}}_{ord}^p(\mathcal{U}, \mathcal{F})
=
\prod\nolimits_{(i_0, \ldots, i_p) \in I^{p + 1}, i_0 < \ldots < i_p}
\mathcal{F}(U_{i_0\ldots i_p}).
$$
This is an abelian group. For
$s \in \check{\mathcal{C}}_{ord}^p(\mathcal{U}, \mathcal{F})$ we denote
$s_{i_0\ldots i_p}$ its value in $\mathcal{F}(U_{i_0\ldots i_p})$.
We define
$$
d : \check{\mathcal{C}}_{ord}^p(\mathcal{U}, \mathcal{F})
\longrightarrow
\check{\mathcal{C}}_{ord}^{p + 1}(\mathcal{U}, \mathcal{F})
$$
by the formula
$$
d(s)_{i_0\ldots i_{p + 1}}
=
\sum\nolimits_{j = 0}^{p + 1}
(-1)^j
s_{i_0\ldots \hat i_j \ldots i_{p + 1}}|_{U_{i_0\ldots i_{p + 1}}}
$$
for any $i_0 < \ldots < i_{p + 1}$. Note that this formula is identical
to Equation (\ref{equation-d-cech}).
It is straightforward to see that $d \circ d = 0$. In other words
$\check{\mathcal{C}}_{ord}^\bullet(\mathcal{U}, \mathcal{F})$ is a complex.

\begin{definition}
\label{definition-ordered-cech-complex}
Let $X$ be a topological space.
Let $\mathcal{U} : U = \bigcup_{i \in I} U_i$ be an open covering.
Assume given a total ordering on $I$.
Let $\mathcal{F}$ be an abelian presheaf on $X$.
The complex $\check{\mathcal{C}}_{ord}^\bullet(\mathcal{U}, \mathcal{F})$
is the {\it ordered {\v C}ech complex} associated to $\mathcal{F}$, the
open covering $\mathcal{U}$ and the given total ordering on $I$.
\end{definition}

\noindent
This complex is sometimes called the alternating {\v C}ech complex.
The reason is that there is an obvious comparison map between
the ordered {\v C}ech complex and the alternating {\v C}ech complex.
Namely, consider the map
$$
c :
\check{\mathcal{C}}_{ord}^\bullet(\mathcal{U}, \mathcal{F})
\longrightarrow
\check{\mathcal{C}}^\bullet(\mathcal{U}, \mathcal{F})
$$
given by the rule
$$
c(s)_{i_0\ldots i_p} =
\left\{
\begin{matrix}
0 &
\text{if} &
i_n = i_m \text{ for some } n \not = m\\
\text{sgn}(\sigma) s_{i_{\sigma(0)}\ldots i_{\sigma(p)}} &
\text{if} &
i_{\sigma(0)} < i_{\sigma(1)} < \ldots < i_{\sigma(p)}
\end{matrix}
\right.
$$
Here $\sigma$ denotes a permutation of $\{0, \ldots, p\}$ and
$\text{sgn}(\sigma)$ denotes its sign. The alternating and ordered
{\v C}ech complexes are often identified in the literature via the map
$c$. Namely we have the following easy lemma.

\begin{lemma}
\label{lemma-ordered-alternating}
Let $X$ be a topological space.
Let $\mathcal{U} : U = \bigcup_{i \in I} U_i$ be an open covering.
Assume $I$ comes equipped with a total ordering.
The map $c$ is a morphism of complexes. In fact it induces
an isomorphism
$$
c : \check{\mathcal{C}}_{ord}^\bullet(\mathcal{U}, \mathcal{F})
\to \check{\mathcal{C}}_{alt}^\bullet(\mathcal{U}, \mathcal{F})
$$
of complexes.
\end{lemma}

\begin{proof}
Omitted.
\end{proof}

\noindent
There is also a map
$$
\pi :
\check{\mathcal{C}}^\bullet(\mathcal{U}, \mathcal{F})
\longrightarrow
\check{\mathcal{C}}_{ord}^\bullet(\mathcal{U}, \mathcal{F})
$$
which is described by the rule
$$
\pi(s)_{i_0\ldots i_p} = s_{i_0\ldots i_p}
$$
whenever $i_0 < i_1 < \ldots < i_p$.

\begin{lemma}
\label{lemma-project-to-ordered}
Let $X$ be a topological space.
Let $\mathcal{U} : U = \bigcup_{i \in I} U_i$ be an open covering.
Assume $I$ comes equipped with a total ordering.
The map $\pi : \check{\mathcal{C}}^\bullet(\mathcal{U}, \mathcal{F})
\to \check{\mathcal{C}}_{ord}^\bullet(\mathcal{U}, \mathcal{F})$
is a morphism of complexes. It induces an isomorphism
$$
\pi : \check{\mathcal{C}}_{alt}^\bullet(\mathcal{U}, \mathcal{F})
\to \check{\mathcal{C}}_{ord}^\bullet(\mathcal{U}, \mathcal{F})
$$
of complexes which is a left inverse to the morphism $c$.
\end{lemma}

\begin{proof}
Omitted.
\end{proof}

\begin{remark}
\label{remark-compared-ordered-complexes}
This means that if we have two total orderings $<_1$ and $<_2$ on
the index set $I$, then we get an isomorphism of complexes
$\tau = \pi_2 \circ c_1 :
\check{\mathcal{C}}_{ord\text{-}1}(\mathcal{U}, \mathcal{F}) \to
\check{\mathcal{C}}_{ord\text{-}2}(\mathcal{U}, \mathcal{F})$.
It is clear that
$$
\tau(s)_{i_0 \ldots i_p} =
\text{sign}(\sigma) s_{i_{\sigma(0)} \ldots i_{\sigma(p)}}
$$
where $i_0 <_1 i_1 <_1 \ldots <_1 i_p$ and
$i_{\sigma(0)} <_2 i_{\sigma(1)} <_2 \ldots <_2 i_{\sigma(p)}$.
This is the sense in which the ordered {\v C}ech complex is independent
of the chosen total ordering.
\end{remark}

\begin{lemma}
\label{lemma-alternating-usual}
Let $X$ be a topological space.
Let $\mathcal{U} : U = \bigcup_{i \in I} U_i$ be an open covering.
Assume $I$ comes equipped with a total ordering.
The map $c \circ \pi$ is homotopic to the identity on
$\check{\mathcal{C}}^\bullet(\mathcal{U}, \mathcal{F})$.
In particular the inclusion map
$\check{\mathcal{C}}_{alt}^\bullet(\mathcal{U}, \mathcal{F}) \to
\check{\mathcal{C}}^\bullet(\mathcal{U}, \mathcal{F})$
is a homotopy equivalence.
\end{lemma}

\begin{proof}
For any multi-index $(i_0, \ldots, i_p) \in I^{p + 1}$ there exists
a unique permutation $\sigma : \{0, \ldots, p\} \to \{0, \ldots, p\}$
such that
$$
i_{\sigma(0)} \leq i_{\sigma(1)} \leq \ldots \leq i_{\sigma(p)}
\quad
\text{and}
\quad
\sigma(j) < \sigma(j + 1)
\quad
\text{if}
\quad
i_{\sigma(j)} = i_{\sigma(j + 1)}.
$$
We denote this permutation $\sigma = \sigma^{i_0 \ldots i_p}$.

\medskip\noindent
For any permutation $\sigma : \{0, \ldots, p\} \to \{0, \ldots, p\}$
and any $a$, $0 \leq a \leq p$ we denote $\sigma_a$
the unique permutation of $\{0, \ldots, p\}$ such that
$\sigma_a(j) = \sigma(j)$ for $0 \leq j < a$ and such that
$\sigma_a(a) < \sigma_a(a + 1) < \ldots < \sigma_a(p)$.
So if $p = 3$ and $\sigma$, $\tau$ are given by
$$
\begin{matrix}
\text{id} & 0 & 1 & 2 & 3 \\
\sigma & 3 & 2 & 1 & 0
\end{matrix}
\quad \text{and} \quad
\begin{matrix}
\text{id} & 0 & 1 & 2 & 3 \\
\tau & 3 & 0 & 2 & 1
\end{matrix}
$$
then we have
$$
\begin{matrix}
\text{id} & 0 & 1 & 2 & 3 \\
\sigma_0 & 0 & 1 & 2 & 3 \\
\sigma_1 & 3 & 0 & 1 & 2 \\
\sigma_2 & 3 & 2 & 0 & 1 \\
\sigma_3 & 3 & 2 & 1 & 0 \\
\end{matrix}
\quad \text{and} \quad
\begin{matrix}
\text{id} & 0 & 1 & 2 & 3 \\
\tau_0 & 0 & 1 & 2 & 3 \\
\tau_1 & 3 & 0 & 1 & 2 \\
\tau_2 & 3 & 0 & 1 & 2 \\
\tau_3 & 3 & 0 & 2 & 1 \\
\end{matrix}
$$
It is clear that always $\sigma_0 = \text{id}$ and $\sigma_p = \sigma$.

\medskip\noindent
Having introduced this notation we define for
$s \in \check{\mathcal{C}}^{p + 1}(\mathcal{U}, \mathcal{F})$
the element $h(s) \in \check{\mathcal{C}}^p(\mathcal{U}, \mathcal{F})$
to be the element with components
\begin{equation}
\label{equation-first-homotopy}
h(s)_{i_0\ldots i_p} =
\sum\nolimits_{0 \leq a \leq p}
(-1)^a \text{sign}(\sigma_a)
s_{i_{\sigma(0)} \ldots i_{\sigma(a)} i_{\sigma_a(a)} \ldots i_{\sigma_a(p)}}
\end{equation}
where $\sigma = \sigma^{i_0 \ldots i_p}$. The index
$i_{\sigma(a)}$ occurs twice in
$i_{\sigma(0)} \ldots i_{\sigma(a)} i_{\sigma_a(a)} \ldots i_{\sigma_a(p)}$
once in the first group of $a + 1$ indices and once in the second group
of $p - a + 1$ indices since $\sigma_a(j) = \sigma(a)$ for some
$j \geq a$ by definition of $\sigma_a$. Hence the sum makes sense since each
of the elements
$s_{i_{\sigma(0)} \ldots i_{\sigma(a)} i_{\sigma_a(a)} \ldots i_{\sigma_a(p)}}$
is defined over the open $U_{i_0 \ldots i_p}$.
Note also that for $a = 0$ we get $s_{i_0 \ldots i_p}$ and
for $a = p$ we get
$(-1)^p \text{sign}(\sigma) s_{i_{\sigma(0)} \ldots i_{\sigma(p)}}$.

\medskip\noindent
We claim that
$$
(dh + hd)(s)_{i_0 \ldots i_p} =
s_{i_0 \ldots i_p} -
\text{sign}(\sigma) s_{i_{\sigma(0)} \ldots i_{\sigma(p)}}
$$
where $\sigma = \sigma^{i_0 \ldots i_p}$. We omit the verification
of this claim. (There is a PARI/gp script called first-homotopy.gp
in the stacks-project subdirectory scripts which can be used to check
finitely many instances of this claim.
We wrote this script to make sure the signs are correct.)
Write
$$
\kappa :
\check{\mathcal{C}}^\bullet(\mathcal{U}, \mathcal{F})
\longrightarrow
\check{\mathcal{C}}^\bullet(\mathcal{U}, \mathcal{F})
$$
for the operator given by the rule
$$
\kappa(s)_{i_0 \ldots i_p} =
\text{sign}(\sigma^{i_0 \ldots i_p}) s_{i_{\sigma(0)} \ldots i_{\sigma(p)}}.
$$
The claim above implies that $\kappa$ is a morphism of complexes and that
$\kappa$ is homotopic to the identity map of the {\v C}ech complex.
This does not immediately imply the lemma since
the image of the operator $\kappa$ is not the alternating subcomplex.
Namely, the image of $\kappa$ is the ``semi-alternating'' complex
$\check{\mathcal{C}}_{semi\text{-}alt}^p(\mathcal{U}, \mathcal{F})$
where $s$ is a $p$-cochain of this complex if and only if
$$
s_{i_0 \ldots i_p} = \text{sign}(\sigma) s_{i_{\sigma(0)} \ldots i_{\sigma(p)}}
$$
for any $(i_0, \ldots, i_p) \in I^{p + 1}$ with
$\sigma = \sigma^{i_0 \ldots i_p}$.
We introduce yet another variant {\v C}ech complex, namely the semi-ordered
{\v C}ech complex defined by
$$
\check{\mathcal{C}}_{semi\text{-}ord}^p(\mathcal{U}, \mathcal{F})
=
\prod\nolimits_{i_0 \leq i_1 \leq \ldots \leq i_p}
\mathcal{F}(U_{i_0 \ldots i_p})
$$
It is easy to see that Equation (\ref{equation-d-cech}) also defines
a differential and hence that we get a complex. It is also clear
(analogous to Lemma \ref{lemma-project-to-ordered}) that the projection map
$$
\check{\mathcal{C}}_{semi\text{-}alt}^\bullet(\mathcal{U}, \mathcal{F})
\longrightarrow
\check{\mathcal{C}}_{semi\text{-}ord}^\bullet(\mathcal{U}, \mathcal{F})
$$
is an isomorphism of complexes.

\medskip\noindent
Hence the Lemma follows if we can show that the obvious inclusion map
$$
\check{\mathcal{C}}_{ord}^p(\mathcal{U}, \mathcal{F})
\longrightarrow
\check{\mathcal{C}}_{semi\text{-}ord}^p(\mathcal{U}, \mathcal{F})
$$
is a homotopy equivalence. To see this we use the homotopy
\begin{equation}
\label{equation-second-homotopy}
h(s)_{i_0 \ldots i_p} =
\left\{
\begin{matrix}
0 & \text{if} & i_0 < i_1 < \ldots < i_p \\
(-1)^a s_{i_0 \ldots i_{a - 1} i_a i_a i_{a + 1} \ldots i_p}
& \text{if} & i_0 < i_1 < \ldots < i_{a - 1} < i_a = i_{a + 1}
\end{matrix}
\right.
\end{equation}
We claim that
$$
(dh + hd)(s)_{i_0 \ldots i_p} =
\left\{
\begin{matrix}
0 & \text{if} & i_0 < i_1 < \ldots < i_p \\
s_{i_0 \ldots i_p}
& \text{else} &
\end{matrix}
\right.
$$
We omit the verification. (There is a PARI/gp script called second-homotopy.gp
in the stacks-project subdirectory scripts which can be used to check
finitely many instances of this claim.
We wrote this script to make sure the signs are correct.)
The claim clearly shows that the composition
$$
\check{\mathcal{C}}_{semi\text{-}ord}^\bullet(\mathcal{U}, \mathcal{F})
\longrightarrow
\check{\mathcal{C}}_{ord}^\bullet(\mathcal{U}, \mathcal{F})
\longrightarrow
\check{\mathcal{C}}_{semi\text{-}ord}^\bullet(\mathcal{U}, \mathcal{F})
$$
of the projection with the natural inclusion
is homotopic to the identity map as desired.
\end{proof}

\begin{lemma}
\label{lemma-alternating-cech-trivial}
Let $X$ be a topological space. Let $\mathcal{F}$ be an abelian presheaf on $X$.
Let $\mathcal{U} : U = \bigcup_{i \in I} U_i$ be an open covering. If
$U_i = U$ for some $i \in I$, then the extended alternating {\v C}ech complex
$$
\mathcal{F}(U) \to \check{\mathcal{C}}_{alt}^\bullet(\mathcal{U}, \mathcal{F})
$$
obtained by putting $\mathcal{F}(U)$ in degree $-1$ with differential given by
the canonical map of $\mathcal{F}(U)$ into
$\check{\mathcal{C}}^0(\mathcal{U}, \mathcal{F})$
is homotopy equivalent to $0$. Similarly, for any total ordering on $I$
the extended ordered {\v C}ech complex
$$
\mathcal{F}(U) \to
\check{\mathcal{C}}_{ord}^\bullet(\mathcal{U}, \mathcal{F})
$$
is homotopy equivalent to $0$.
\end{lemma}

\begin{proof}[First proof]
Combine Lemmas \ref{lemma-cech-trivial} and \ref{lemma-alternating-usual}.
\end{proof}

\begin{proof}[Second proof]
Since the alternating and ordered {\v C}ech complexes are isomorphic
it suffices to prove this for the ordered one.
We will use standard notation: a cochain $s$ of degree $p$
in the extended ordered {\v C}ech complex has the form
$s = (s_{i_0 \ldots i_p})$ where $s_{i_0 \ldots i_p}$ is in
$\mathcal{F}(U_{i_0 \ldots i_p})$ and $i_0 < \ldots < i_p$.
With this notation we have
$$
d(x)_{i_0 \ldots i_{p + 1}} =
\sum\nolimits_j (-1)^j x_{i_0 \ldots \hat i_j \ldots i_p}
$$
Fix an index $i \in I$ with $U = U_i$.
As homotopy we use the maps
$$
h : \text{cochains of degree }p + 1 \to \text{cochains of degree }p
$$
given by the rule
$$
h(s)_{i_0 \ldots i_p} = 0 \text{ if } i \in \{i_0, \ldots, i_p\}
\text{ and }
h(s)_{i_0 \ldots i_p} = 
(-1)^j s_{i_0 \ldots i_j i i_{j + 1} \ldots i_p} \text{ if not}
$$
Here $j$ is the unique index such that $i_j < i < i_{j + 1}$ in the
second case; also, since $U = U_i$ we have the equality
$$
\mathcal{F}(U_{i_0 \ldots i_p}) =
\mathcal{F}(U_{i_0 \ldots i_j i i_{j + 1} \ldots i_p})
$$
which we can use to make sense of thinking of
$(-1)^j s_{i_0 \ldots i_j i i_{j + 1} \ldots i_p}$
as an element of $\mathcal{F}(U_{i_0 \ldots i_p})$.
We will show by a computation that $d h + h d$ equals
the negative of the identity map which finishes the proof.
To do this fix $s$ a cochain of degree $p$ and let
$i_0 < \ldots < i_p$ be elements of $I$.

\medskip\noindent
Case I: $i \in \{i_0, \ldots, i_p\}$. Say $i = i_t$. Then we have
$h(d(s))_{i_0 \ldots i_p} = 0$. On the other hand we have
$$
d(h(s))_{i_0 \ldots i_p} =
\sum (-1)^j h(s)_{i_0 \ldots \hat i_j \ldots i_p} =
(-1)^t h(s)_{i_0 \ldots \hat i \ldots i_p} =
(-1)^t (-1)^{t - 1} s_{i_0 \ldots i_p}
$$
Thus $(dh + hd)(s)_{i_0 \ldots i_p} = -s_{i_0 \ldots i_p}$ as desired.

\medskip\noindent
Case II: $i \not \in \{i_0, \ldots, i_p\}$. Let $j$ be such that
$i_j < i < i_{j + 1}$. Then we see that
\begin{align*}
h(d(s))_{i_0 \ldots i_p}
& =
(-1)^j d(s)_{i_0 \ldots i_j i i_{j + 1} \ldots i_p} \\
& =
\sum\nolimits_{j' \leq j} (-1)^{j + j'}
s_{i_0 \ldots \hat i_{j'} \ldots i_j i i_{j + 1} \ldots i_p} -
s_{i_0 \ldots i_p} \\
&
+ \sum\nolimits_{j' > j} (-1)^{j + j' + 1}
s_{i_0 \ldots i_j i i_{j + 1} \ldots \hat i_{j'} \ldots i_p}
\end{align*}
On the other hand we have
\begin{align*}
d(h(s))_{i_0 \ldots i_p}
& =
\sum\nolimits_{j'} (-1)^{j'} h(s)_{i_0 \ldots \hat i_{j'} \ldots i_p} \\
& =
\sum\nolimits_{j' \leq j} (-1)^{j' + j - 1}
s_{i_0 \ldots \hat i_{j'} \ldots i_j i i_{j + 1} \ldots i_p} \\
& +
\sum\nolimits_{j' > j} (-1)^{j' + j}
s_{i_0 \ldots i_j i i_{j + 1} \ldots \hat i_{j'} \ldots i_p}
\end{align*}
Adding these up we obtain
$(dh + hd)(s)_{i_0 \ldots i_p} = - s_{i_0 \ldots i_p}$
as desired.
\end{proof}





